\documentclass[10pt,oneside]{article}

%%%%%%%%%%%%%
\setlength{\textheight}{8.75in} %Letter is 11in, less 2 for margins, less 0.25 for footer
\setlength{\oddsidemargin}{0.0in} %gets +1inc
\setlength{\evensidemargin}{0.0in} %gets +1inch
\setlength{\textwidth}{6.50in} %Letter is 8.5, less 2 inches for margins
\setlength{\topmargin}{0.5in}
\setlength{\headheight}{0in}
\setlength{\headsep}{0in}
\setlength{\parindent}{0.25in}
%%%%%%%%%%%%

\usepackage[numbers]{natbib}
\usepackage{graphicx}
\usepackage[colorlinks=true,citecolor=black,urlcolor=blue]{hyperref}

\title{%Hack to get the logo on the PDF front page:
\vspace{-1.5in}
\includegraphics[width=0.5\textwidth]{biopython.jpg}\\
%Hack to get some white space using a blank line:
~\\Biopython Project Update 2016}
\author{
	\underline{Christian Brueffer}\footnote{Division of Oncology and Pathology, Department of Clinical Sciences, Lund University, Lund, SE. Email: \href{mailto:christian.brueffer@med.lu.se}{christian.brueffer@med.lu.se}},
    Jo\~{a}o Rodrigues\footnote{Bijvoet Center for Biomolecular Research, Utrecht University, Utrecht, NL},
    Tiago Ant\~{a}o \footnote{Vector Biology Department, Liverpool School of Tropical Medicine, Pembroke Place, UK},
    Peter Cock\footnote{Information and Computational Sciences, James Hutton Institute (formerly SCRI), Invergowrie, Dundee, UK},
    Eric Talevich\footnote{Department of Dermatology, University of California San Francisco, San Francisco, CA, USA},
    Michiel de Hoon\footnote{Division of Genomic Technologies, RIKEN Center for Life Science Technologies, Yokohama, JP},
		\\
    Wibowo Arindrarto\footnote{Sequencing Analysis Support Core, Leiden University Medical Center, Leiden, NL},
    and the Biopython Contributors}
\date{17\textsuperscript{th} Bioinformatics Open Source Conference (BOSC) 2016, Orlando, USA}

\begin{document}
%\pagestyle{empty}
%\thispagestyle{empty}
\maketitle
%\pagestyle{empty}
\thispagestyle{empty}

\vspace{-0.2in}
\noindent
Website: \url{http://biopython.org} \\
Repository: \url{https://github.com/biopython/biopython} \\
License: Biopython License Agreement (MIT style, see \url{http://www.biopython.org/DIST/LICENSE}) \\

The Biopython Project is a long-running distributed collaborative effort,
supported by the Open Bioinformatics Foundation, which develops a freely
available Python library for biological computation \cite{AppNote}.

We present here details of the latest Biopython release - version 1.66. New
features include: extended Bio.KEGG and Bio.Graphics modules to support drawing
KEGG pathways with transparency; extended the ``abi'' Bio.SeqIO parser to decode
almost all documented fields used by ABIF instruments; added a QCPSuperimposer
module using the Quaternion Characteristic Polynomial algorithm for superimposing
structures to Bio.PDB; and extended the Bio.Entrez module to implement the NCBI
Entrez Citation Matching function and to support NCBI XML files with XSD schemas.
Additionally we fixed miscellaneous bugs, enhanced our test suite and continued our
efforts to abide by the PEP8 coding style guidelines.

We are currently preparing a new release - version 1.67 - that will deprecate the
ability to compare SeqRecord objects with ``=='', which sometimes lead to surprising
results.  In addition it will feature updates to Bio.Data to include NCBI genetic
code table 25, Candidate Division SR1 and Gracilibacteria; and updates to BioSQL to
use foreign keys with SQLite3 databases.

Our continuous integration process on GitHub has been enhanced by including
external services services like Landscape, Quantified Code and Codecov to perform
quality review, test coverage analysis and generation of quality metrics.

Finally, our range of Docker containers has been greatly enhanced. In addition to
a basic containers that include Python 2 and 3 with Biopython and all its
dependencies, as well as BioSQL container, we now also two versions of Jupyter
notebook containers: a basic one, and a version including the Biopython tutorial as
notebooks.


\begin{thebibliography}{}

\bibitem[Cock {\it et al}., 2009]{AppNote}Cock, P.J.A., Antao, T., Chang, J.T., Chapman, B.A., Cox, C.J., Dalke, A., Friedberg, I., Hamelryck, T., Kauff, F., Wilczynski, B., de Hoon, M.J. (2009) Biopython: freely available Python tools for computational molecular biology and bioinformatics. {\it Bioinformatics} {\bf 25}(11) 1422-3. \href{http://dx.doi.org/10.1093/bioinformatics/btp163}{doi:10.1093/bioinformatics/btp163}

\end{thebibliography}

\end{document}
